\section{Acordo de Não Divulgação (NDA)}

\paragraph{}
Antes do início das atividades de teste de penetração, deverá ser firmado um
Acordo de Não Divulgação (\textit{Non-Disclosure Agreement -- NDA}) entre a
equipe responsável pela execução dos testes e os responsáveis pelo Laboratório
de Segurança. O objetivo do acordo é proteger todas as informações obtidas
durante o planejamento, execução e encerramento do Pentest.

\paragraph{}
Serão consideradas informações confidenciais todos os dados técnicos,
administrativos e operacionais acessados durante o projeto, incluindo, mas não
se limitando a:

\begin{itemize}
\item Vulnerabilidades identificadas;
\item Evidências de exploração;
\item Relatórios técnicos e executivos;
\item Credenciais eventualmente disponibilizadas ou encontradas durante os
testes;
\item Informações sobre infraestrutura, servidores, aplicações e serviços;
\item Configurações de segurança;
\item Registros de comunicação entre as partes envolvidas.
\end{itemize}

\paragraph{}
As informações obtidas durante o Pentest não poderão ser divulgadas,
compartilhadas ou utilizadas para qualquer finalidade diferente daquela
prevista no escopo do projeto sem autorização formal dos responsáveis pelo
Laboratório de Segurança.

\paragraph{}
Todos os participantes envolvidos nas atividades deverão assinar o NDA antes do
início dos testes, comprometendo-se a preservar o sigilo das informações
obtidas e a adotar medidas adequadas para sua proteção.

\paragraph{}
O acordo deverá estabelecer claramente as obrigações das partes, os limites de
uso das informações, os procedimentos de armazenamento e descarte seguro dos
dados coletados e as penalidades aplicáveis em caso de violação da
confidencialidade.

\paragraph{}
Também deverão ser observados os requisitos legais e regulatórios aplicáveis,
incluindo normas relacionadas à proteção de dados e privacidade das informações
eventualmente tratadas durante a execução do Pentest.

\paragraph{}
O compromisso de confidencialidade permanecerá válido mesmo após a conclusão do
projeto, abrangendo todos os relatórios, evidências e informações obtidas
durante a realização dos testes de segurança.

\paragraph{}
Como condição para o início das atividades, deverá ser confirmada a assinatura
do NDA por todas as partes envolvidas e a aprovação formal da documentação do
projeto, incluindo escopo, plano de testes, responsabilidades, regras de
comunicação e critérios de interrupção das atividades.
